\chapter*{Préface}
\addcontentsline{toc}{chapter}{Préface}
Ce livre raconte une histoire mathématique récente, telle qu'elle s'est réellement déroulée : une suite de conjectures,
de démonstrations, de vérifications par ordinateur et parfois de réfutations. Il s'adresse aux élèves des classes
préparatoires (MPSI, MP2I, MP, MPI) et des écoles d'ingénieurs. Le point de départ est un théorème d'openai/math, dont la démonstration, produite par une IA et non relue par des pairs, est vérifiée par le noyau de Lean~4 : le réseau triangulaire minimise l'énergie de toute configuration de points du plan pour une large classe
d'interactions répulsives. Le point d'arrivée est une série de questions ouvertes, en physique quantique et en
génétique des populations.

Chaque résultat est présenté trois fois.
\begin{itemize}
\item D'abord comme une \emph{démonstration humaine}, pas à pas, au niveau des classes préparatoires.
\item Ensuite comme un énoncé \emph{Lean~4} vérifié par le noyau.
\item Enfin comme un calcul numérique avec \emph{rusty-SUNDIALS} (solveurs CVODE en Rust) et \emph{LeanFlow}
(certificats de données vérifiés par Lean), illustré par des jeux de données publics.
\end{itemize}

Une seule règle guide la lecture : chaque affirmation porte un \emph{niveau de confiance}, selon la méthode Elenchus.
\begin{itemize}
\item \tierA{} : vérifié par le noyau Lean. Cela suppose trois vérifications : la liste des axiomes se limite à
\lean{propext}, \lean{Classical.choice} et \lean{Quot.sound}, le noyau rejoue la preuve (\lean{leanchecker}), et la
recherche des constructions qui contournent le noyau ne trouve rien. La liste des axiomes seule peut être trompée ;
le chapitre~1 explique comment.
\item \tierAs{} : variante de \tierA{} quand le rejeu du noyau n'a pas (encore) été fait pour le module concerné ; la liste des axiomes, les témoins plantés et la recompilation indépendante restent exigés.
\item \tierL{} : résultat publié, lu et cité précisément.
\item \tierC{} : argument humain, non vérifié par machine.
\item \tierX{} : calcul en virgule flottante. Il oriente la recherche mais ne démontre rien.
\end{itemize}
Lisez le niveau avant la phrase.

\paragraph{Honnêteté.}
\begin{itemize}
\item Les résultats de ce livre ont été produits par des agents d'IA (Claude) sous la direction de l'auteur.
\item Chaque article a été audité par un agent d'IA distinct ; aucun mathématicien humain ne les a encore relus.
\item Plusieurs résultats dépendent de théorèmes de la publication openai/math, eux aussi produits par une IA et non
encore évalués par des pairs.
\item Plusieurs chapitres rapportent des \emph{résultats négatifs} : une conjecture réfutée, une méthode qui ne s'étend
pas, une transition qui n'a pas lieu. Ils sont aussi instructifs que les autres.
\item Aucune conjecture n'est déclarée démontrée au-delà de son niveau. En particulier, la conjecture de
Sandier--Serfaty n'est \emph{pas} démontrée ici : le chapitre~2 explique exactement ce qui a été vérifié et ce qui
reste à faire.
\end{itemize}

\paragraph{Plan.}
\begin{itemize}
\item Le chapitre~1 pose les fondations : réseaux, fonctions complètement monotones et vérification formelle.
\item Le chapitre~2 suit le fil des formalisations, de l'optimalité universelle aux potentiels de la physique et à la
conjecture de Sandier--Serfaty.
\item Le chapitre~3 présente une surprise quantique : pour certaines interactions, le réseau triangulaire n'est plus
optimal dès qu'on tient compte de l'énergie de point zéro.
\item Le chapitre~4 vérifie en Lean des théorèmes classiques de génétique des populations, jusqu'au théorème
isotherme des graphes évolutifs.
\item Le chapitre~5 réunit cinq cas d'usage avec données, codes et certificats.
\end{itemize}
Chaque chapitre se termine par des exercices.

\paragraph{Pour aller plus loin.} Tous les articles cités ont un DOI Zenodo. Le code est public : AutoevolveAI,
rusty-SUNDIALS, LeanFlow, SocrateAI-Scientific-Genetic et SocrateAI-Scientific-Agora-LeanMaster.
